% Project Journey Portfolio — group submission (one PDF).
% Build: latexmk -pdf project-journey-portfolio.tex
\documentclass[11pt,a4paper]{article}
% Shared style for the StudyOS Thesis-Finder course submissions.
% Included by every submission document via % Shared style for the StudyOS Thesis-Finder course submissions.
% Included by every submission document via % Shared style for the StudyOS Thesis-Finder course submissions.
% Included by every submission document via \input{preamble}.
% Plain black-and-white house style: no colour, prints identically in mono.

\usepackage[T1]{fontenc}
\usepackage[utf8]{inputenc}
\usepackage{lmodern}
\usepackage[english]{babel}
\usepackage{microtype}
\usepackage[a4paper,margin=2.4cm]{geometry}
\usepackage{xcolor}
\usepackage{titlesec}
\usepackage{enumitem}
\usepackage{booktabs}
\usepackage{tabularx}
\usepackage{tcolorbox}
\tcbuselibrary{breakable,skins}
\usepackage{pifont}
\usepackage{xurl}
\usepackage[hidelinks]{hyperref}

\urlstyle{tt}

% --- Palette (plain black and white) ----------------------------------------
% Everything is black on white; the only greys are for a hairline rule and the
% callout title bar, so the documents print identically in mono.
\definecolor{sosAccent}{gray}{0}
\definecolor{sosDark}{gray}{0}
\definecolor{sosOpen}{gray}{0}
\definecolor{sosOpenBg}{gray}{1}
\definecolor{sosRule}{gray}{0.6}

% --- Headings ---------------------------------------------------------------
\titleformat{\section}{\normalfont\large\bfseries}{\thesection}{0.6em}{}
\titleformat{\subsection}{\normalfont\normalsize\bfseries\itshape}{\thesubsection}{0.6em}{}
\titlespacing*{\section}{0pt}{1.1\baselineskip}{0.35\baselineskip}
\titlespacing*{\subsection}{0pt}{0.7\baselineskip}{0.2\baselineskip}

\setlist[itemize]{leftmargin=1.2em,itemsep=0.15em,topsep=0.25em,parsep=0pt}
\setlist[enumerate]{leftmargin=1.4em,itemsep=0.15em,topsep=0.25em,parsep=0pt}

% --- Semantic helpers -------------------------------------------------------
% Inline reference to a file or path inside the project repository.
\newcommand{\repofile}[1]{\path{#1}}

% One "the assignment asks for ..." lead-in line.
\newcommand{\asks}[1]{\noindent\textbf{Assignment asks for:} #1\par\vspace{0.3em}}

% Open question the team still has to answer (was <mark>highlighted</mark> in the
% markdown draft).
\newenvironment{openbox}[1][Open --- needs the team]{%
  \begin{tcolorbox}[
    breakable,
    colback=white,
    colframe=black,
    colbacktitle=black!8,
    coltitle=black,
    boxrule=0.4pt,
    titlerule=0pt,
    left=0.6em,right=0.6em,top=0.5em,bottom=0.5em,
    title={#1},
    fonttitle=\bfseries\small,
    sharp corners,
  ]\small
}{%
  \end{tcolorbox}
}

\NewDocumentEnvironment{infobox}
  { O{Open --- needs the team} O{} }
{%
  \begin{tcolorbox}[
    breakable,
    colback=white,
    colframe=black,
    colbacktitle=black!8,
    coltitle=black,
    boxrule=0.4pt,
    titlerule=0pt,
    left=0.6em,
    right=0.6em,
    top=0.5em,
    bottom=0.5em,
    title={#1},
    fonttitle=\bfseries\small,
    sharp corners,
    #2
  ]\small
}{%
  \end{tcolorbox}
}


% Status markers for the submission checklist.
\newcommand{\stdone}{\ding{51}}
\newcommand{\stpartial}{\ding{108}}
\newcommand{\stgap}{\ding{55}}

% Title block.
\newcommand{\sostitle}[3]{%
  \begin{center}
    {\LARGE\bfseries #1}\\[0.35em]
    {\large\itshape #2}\\[0.5em]
    {\small #3}
  \end{center}
  \vspace{0.4em}
  {\color{sosRule}\hrule height 0.4pt}
  \vspace{0.9em}
}
.
% Plain black-and-white house style: no colour, prints identically in mono.

\usepackage[T1]{fontenc}
\usepackage[utf8]{inputenc}
\usepackage{lmodern}
\usepackage[english]{babel}
\usepackage{microtype}
\usepackage[a4paper,margin=2.4cm]{geometry}
\usepackage{xcolor}
\usepackage{titlesec}
\usepackage{enumitem}
\usepackage{booktabs}
\usepackage{tabularx}
\usepackage{tcolorbox}
\tcbuselibrary{breakable,skins}
\usepackage{pifont}
\usepackage{xurl}
\usepackage[hidelinks]{hyperref}

\urlstyle{tt}

% --- Palette (plain black and white) ----------------------------------------
% Everything is black on white; the only greys are for a hairline rule and the
% callout title bar, so the documents print identically in mono.
\definecolor{sosAccent}{gray}{0}
\definecolor{sosDark}{gray}{0}
\definecolor{sosOpen}{gray}{0}
\definecolor{sosOpenBg}{gray}{1}
\definecolor{sosRule}{gray}{0.6}

% --- Headings ---------------------------------------------------------------
\titleformat{\section}{\normalfont\large\bfseries}{\thesection}{0.6em}{}
\titleformat{\subsection}{\normalfont\normalsize\bfseries\itshape}{\thesubsection}{0.6em}{}
\titlespacing*{\section}{0pt}{1.1\baselineskip}{0.35\baselineskip}
\titlespacing*{\subsection}{0pt}{0.7\baselineskip}{0.2\baselineskip}

\setlist[itemize]{leftmargin=1.2em,itemsep=0.15em,topsep=0.25em,parsep=0pt}
\setlist[enumerate]{leftmargin=1.4em,itemsep=0.15em,topsep=0.25em,parsep=0pt}

% --- Semantic helpers -------------------------------------------------------
% Inline reference to a file or path inside the project repository.
\newcommand{\repofile}[1]{\path{#1}}

% One "the assignment asks for ..." lead-in line.
\newcommand{\asks}[1]{\noindent\textbf{Assignment asks for:} #1\par\vspace{0.3em}}

% Open question the team still has to answer (was <mark>highlighted</mark> in the
% markdown draft).
\newenvironment{openbox}[1][Open --- needs the team]{%
  \begin{tcolorbox}[
    breakable,
    colback=white,
    colframe=black,
    colbacktitle=black!8,
    coltitle=black,
    boxrule=0.4pt,
    titlerule=0pt,
    left=0.6em,right=0.6em,top=0.5em,bottom=0.5em,
    title={#1},
    fonttitle=\bfseries\small,
    sharp corners,
  ]\small
}{%
  \end{tcolorbox}
}

\NewDocumentEnvironment{infobox}
  { O{Open --- needs the team} O{} }
{%
  \begin{tcolorbox}[
    breakable,
    colback=white,
    colframe=black,
    colbacktitle=black!8,
    coltitle=black,
    boxrule=0.4pt,
    titlerule=0pt,
    left=0.6em,
    right=0.6em,
    top=0.5em,
    bottom=0.5em,
    title={#1},
    fonttitle=\bfseries\small,
    sharp corners,
    #2
  ]\small
}{%
  \end{tcolorbox}
}


% Status markers for the submission checklist.
\newcommand{\stdone}{\ding{51}}
\newcommand{\stpartial}{\ding{108}}
\newcommand{\stgap}{\ding{55}}

% Title block.
\newcommand{\sostitle}[3]{%
  \begin{center}
    {\LARGE\bfseries #1}\\[0.35em]
    {\large\itshape #2}\\[0.5em]
    {\small #3}
  \end{center}
  \vspace{0.4em}
  {\color{sosRule}\hrule height 0.4pt}
  \vspace{0.9em}
}
.
% Plain black-and-white house style: no colour, prints identically in mono.

\usepackage[T1]{fontenc}
\usepackage[utf8]{inputenc}
\usepackage{lmodern}
\usepackage[english]{babel}
\usepackage{microtype}
\usepackage[a4paper,margin=2.4cm]{geometry}
\usepackage{xcolor}
\usepackage{titlesec}
\usepackage{enumitem}
\usepackage{booktabs}
\usepackage{tabularx}
\usepackage{tcolorbox}
\tcbuselibrary{breakable,skins}
\usepackage{pifont}
\usepackage{xurl}
\usepackage[hidelinks]{hyperref}

\urlstyle{tt}

% --- Palette (plain black and white) ----------------------------------------
% Everything is black on white; the only greys are for a hairline rule and the
% callout title bar, so the documents print identically in mono.
\definecolor{sosAccent}{gray}{0}
\definecolor{sosDark}{gray}{0}
\definecolor{sosOpen}{gray}{0}
\definecolor{sosOpenBg}{gray}{1}
\definecolor{sosRule}{gray}{0.6}

% --- Headings ---------------------------------------------------------------
\titleformat{\section}{\normalfont\large\bfseries}{\thesection}{0.6em}{}
\titleformat{\subsection}{\normalfont\normalsize\bfseries\itshape}{\thesubsection}{0.6em}{}
\titlespacing*{\section}{0pt}{1.1\baselineskip}{0.35\baselineskip}
\titlespacing*{\subsection}{0pt}{0.7\baselineskip}{0.2\baselineskip}

\setlist[itemize]{leftmargin=1.2em,itemsep=0.15em,topsep=0.25em,parsep=0pt}
\setlist[enumerate]{leftmargin=1.4em,itemsep=0.15em,topsep=0.25em,parsep=0pt}

% --- Semantic helpers -------------------------------------------------------
% Inline reference to a file or path inside the project repository.
\newcommand{\repofile}[1]{\path{#1}}

% One "the assignment asks for ..." lead-in line.
\newcommand{\asks}[1]{\noindent\textbf{Assignment asks for:} #1\par\vspace{0.3em}}

% Open question the team still has to answer (was <mark>highlighted</mark> in the
% markdown draft).
\newenvironment{openbox}[1][Open --- needs the team]{%
  \begin{tcolorbox}[
    breakable,
    colback=white,
    colframe=black,
    colbacktitle=black!8,
    coltitle=black,
    boxrule=0.4pt,
    titlerule=0pt,
    left=0.6em,right=0.6em,top=0.5em,bottom=0.5em,
    title={#1},
    fonttitle=\bfseries\small,
    sharp corners,
  ]\small
}{%
  \end{tcolorbox}
}

\NewDocumentEnvironment{infobox}
  { O{Open --- needs the team} O{} }
{%
  \begin{tcolorbox}[
    breakable,
    colback=white,
    colframe=black,
    colbacktitle=black!8,
    coltitle=black,
    boxrule=0.4pt,
    titlerule=0pt,
    left=0.6em,
    right=0.6em,
    top=0.5em,
    bottom=0.5em,
    title={#1},
    fonttitle=\bfseries\small,
    sharp corners,
    #2
  ]\small
}{%
  \end{tcolorbox}
}


% Status markers for the submission checklist.
\newcommand{\stdone}{\ding{51}}
\newcommand{\stpartial}{\ding{108}}
\newcommand{\stgap}{\ding{55}}

% Title block.
\newcommand{\sostitle}[3]{%
  \begin{center}
    {\LARGE\bfseries #1}\\[0.35em]
    {\large\itshape #2}\\[0.5em]
    {\small #3}
  \end{center}
  \vspace{0.4em}
  {\color{sosRule}\hrule height 0.4pt}
  \vspace{0.9em}
}


\hypersetup{pdftitle={Project Journey Portfolio — StudyOS Thesis-Finder}}

\begin{document}

\sostitle{Project Journey Portfolio}{StudyOS Thesis-Finder --- working draft}{%
  Group submission \quad\textbullet\quad Domenique, Max, Valentin \quad\textbullet\quad University of Tübingen}

\noindent
Info-gathering version. Each section mirrors the assignment structure
(What / Evidence / So what / Now what) and states what we already have from the
repository. Everything marked in an \textbf{Open} box is a
question the team still needs to answer --- these are the gaps between what is
documented and what the assignment explicitly asks for.

\vspace{0.5em}
\begin{tcolorbox}[colback=white,colframe=sosRule,boxrule=0.6pt,left=0.7em,right=0.7em,top=0.5em,bottom=0.5em]
\textbf{The journey logic in one paragraph (draft):} We started with a matching
platform (app + database) for thesis search. Early conversations with professors
broke the core premise: the supply side won't feed a central platform, and a
curated database rots. We reframed the problem from ``matching'' to the
student's cold start, pivoted to a database-free agent skill that searches the
live web, built and self-evaluated it against a plain-Claude baseline, and
shipped it as a public, installable package. What's still missing is testing in
the hands of real students --- the protocol is ready, the test hasn't run yet.
\end{tcolorbox}

% ---------------------------------------------------------------------------
\section{Initial Project Idea}

\asks{initial idea, motivation, initial assumptions, first problem idea, target
group, an early proposal/sketch/user story --- plus why it seemed promising and
what needed checking next.}

\subsection*{What we have}
\begin{itemize}
  \item Initial idea: a hosted web platform that matches students to thesis
    advisors --- student profiles on one side, a database of chairs and thesis
    topics on the other, classic matching-product architecture (backend +
    database + frontend).
  \item Target group: students at our university searching for a thesis,
    starting with computer science.
  \item First problem idea: the start of a thesis search is painful ---
    information is scattered, students don't know their options.
  \item Initial assumptions (implicit at the time, worth stating explicitly
    because they all broke later): (a)~a central platform is the right vehicle,
    (b)~professors/chairs will participate and list topics, (c)~the database can
    be kept fresh.
  \item The app was built to roughly 90\,\% feature-completeness before we
    abandoned it --- that's an honest and interesting data point for the
    portfolio, not something to hide.
\end{itemize}

\subsection*{Evidence to attach}
\begin{itemize}
  \item Pre-pivot architecture/status snapshot:
    \repofile{docs/thesis-report/00-problem-and-research/2026-06-09-project-context-pre-pivot.md}
  \item First entry of the decision log:
    \repofile{docs/thesis-report/decision-log.md}
  \item Early plan and outreach documents: various private conversations ---
    please request them if you need them to rate the submission, but since they
    are private data we cannot release them publicly.
\end{itemize}

\subsection*{So what (draft)}
\begin{itemize}
  \item Promising because the pain is universally recognized among students, and
    ``matching platform'' is a well-understood product shape that felt buildable.
\end{itemize}

\subsection*{Now what (draft)}
\begin{itemize}
  \item The obvious next check: is there actual demand on both sides --- do
    students want this, and will professors participate?
\end{itemize}

\begin{openbox}
\begin{itemize}
  \item What was the actual trigger for this project --- our own thesis-search
    frustration, the course prompt, something else? The motivation has to feel
    real, not retro-fitted.
  \item Does one genuinely early artifact survive (first slide deck, sketch, or
    user story)? The assignment explicitly wants one piece of early evidence ---
    pick the single best one to include.
  \item Can we reconstruct the very first ``we aim to solve X for Y by doing Z''
    sentence as we would have written it on day one? We need it as the baseline
    that Section~3 shows changing.
\end{itemize}
\end{openbox}

% ---------------------------------------------------------------------------
\section{Discovering and Understanding Users}

\asks{user groups discovered and selected, how they were reached, the research
approach, main insights, a persona or user story --- plus what we learned that
the initial idea alone couldn't tell us, and how that shaped the direction.}

\subsection*{What we have}
\begin{itemize}
  \item Two user groups, and this distinction matters for the story: students
    are the demand side, professors/chairs are the supply side a platform would
    depend on.
  \item We interviewed a substantial number of professors. Main documented
    insights:
    \begin{itemize}
      \item Roughly half rejected a central platform outright; only about a
        third were conditionally open to it.
      \item A comparable topic-listing platform at the university had already
        failed --- for incentive reasons, not UI reasons: nobody wants to list
        and maintain open topics.
      \item Individual professors reacted much more positively to a different
        framing than to our original platform idea.
    \end{itemize}
  \item The course professor later gave feedback that shifted what the product
    is even \emph{for}: the real value is that students reflect on what they
    want --- not the concrete proposals the tool outputs. That insight is now
    built into the skill (it records the student's self-understanding before and
    after the search).
\end{itemize}

\subsection*{Evidence to attach}
\begin{itemize}
  \item Interview findings and quotes:
    \repofile{docs/thesis-report/00-problem-and-research/README.md} and the
    professor research package in the same folder.
  \item Inventory of every documented user contact, with quotes:
    \repofile{findings/no_db_universal_skill/2026-08-09-user-study-inventory.md}
\end{itemize}

\subsection*{So what (draft)}
\begin{itemize}
  \item The initial idea could not have told us that the \emph{supply side} is
    the bottleneck. Interviews revealed the platform premise fails regardless of
    how good the product is --- participation and freshness are incentive
    problems, not engineering problems.
\end{itemize}

\subsection*{Now what (draft)}
\begin{itemize}
  \item This forced the reframing in Section~3: solve the student's problem
    without depending on professor participation or on curated data.
\end{itemize}

\begin{openbox}
\begin{itemize}
  \item How were the interviews actually done --- how were professors selected
    and reached, was there a fixed question guide or informal conversations, and
    who conducted them? The repo itself flags that this method was never written
    down.
  \item What student-side input existed before the pivot? Right now the
    documented pivot evidence is professor-side only, but the product targets
    students --- if we talked to fellow students informally (friends,
    Fachschaft, ourselves as users), that must be described; if we didn't, we
    should say so honestly and explain why professor evidence was enough.
  \item Can we write one short persona or user story for our target student? The
    assignment asks for one explicitly and we have none.
  \item What is the single most surprising thing a user said to us? One strong
    quote carries more than a paragraph of summary.
\end{itemize}
\end{openbox}

% ---------------------------------------------------------------------------
\section{Defining the Problem}

\asks{initial vs.\ refined problem definition, which assumptions changed and
why, and the fit between users, context, and problem.}

\subsection*{What we have}
\begin{itemize}
  \item Initial problem definition (implicit): ``students can't find matching
    advisors $\rightarrow$ build a matching platform.''
  \item Refined problem definition: the hardest part is the \textbf{cold start}
    --- thesis information is scattered across chair websites, students don't
    know what they don't know, and any curated database only covers one faculty
    and only stays correct while someone maintains it.
  \item The two questions that triggered the rethink are literally preserved in
    our meeting notes: \emph{``Do we even need a backend?''} and \emph{``Where's
    the difference to just using Claude without a skill?''}
  \item Assumptions that broke, with the evidence that broke them: professor
    participation (interviews), database freshness (the failed predecessor
    platform + our own maintenance backlog), and willingness to use a separate
    external tool (nobody opens yet another app for a one-time task).
\end{itemize}

\subsection*{Evidence to attach}
\begin{itemize}
  \item Raw pivot meeting notes:
    \repofile{docs/thesis-report/01-the-pivot/2026-06-25-besprechung-notes.md}
  \item Reframed problem statement:
    \repofile{docs/thesis-report/01-the-pivot/2026-06-26-vision-no-db.md}
  \item Decision log entry for the pivot:
    \repofile{docs/thesis-report/decision-log.md}
\end{itemize}

\subsection*{So what (draft)}
\begin{itemize}
  \item The refined definition is better because it survives the broken
    assumptions: it doesn't require anyone's participation, doesn't rot, and
    meets students where they already are. It also honestly shrinks the claim
    --- from ``we match you'' to ``we map your options.''
\end{itemize}

\subsection*{Now what (draft)}
\begin{itemize}
  \item The refined problem directly dictates the solution constraints in
    Section~4: no database, no backend, live sources only, all faculties, zero
    maintenance.
\end{itemize}

\begin{openbox}
\begin{itemize}
  \item Write both problem definitions as one clean sentence each (``We aim to
    solve X for Y by doing Z''), initial and refined, so the before/after is
    visible at a glance. Draft them, then agree as a group.
  \item How did the team actually converge on the reframing --- immediate
    agreement or real back-and-forth? One honest sentence about the group
    dynamic belongs here.
\end{itemize}
\end{openbox}

% ---------------------------------------------------------------------------
\section{Proposing a Solution}

\asks{solution concept, value proposition, alternatives considered, feature
prioritization, and the fit between users, problem, and solution.}

\subsection*{What we have}
\begin{itemize}
  \item Solution concept: a portable, database-free \textbf{agent skill package}
    --- markdown instructions that teach an LLM how to interview a student about
    interests/skills/constraints and then search the live web in two passes
    (discover candidates from multiple source axes, then verify and enrich each
    one), producing a \textbf{map of options} with rationale and evidence, not a
    single answer.
  \item Value proposition (draft, needs a group-agreed sentence): always current
    because it reads the live web, covers all faculties immediately, requires
    zero maintenance, and lives inside the tool students already use --- no
    separate app to install or account to create.
  \item Alternatives considered, documented: the original app with database
    (built, abandoned); a database-backed skill vs.\ the database-free skill ---
    genuinely contested inside the team, resolved by pursuing both in parallel
    until testing showed no-DB wins on both user value and maintenance cost.
  \item Why it beats ``just ask Claude'': the skill carries encoded discovery
    rules (source axes, verification checks, interview discipline) --- this was
    reasoned out explicitly and later measured (Section~6).
  \item Feature prioritization, documented: university arm first and proven,
    companies as a second phase on the same principle, other universities and
    Bachelor's theses explicitly out of scope.
  \item The concept was deliberately stress-tested against named risks before we
    committed to it.
\end{itemize}

\subsection*{Evidence to attach}
\begin{itemize}
  \item Full solution vision:
    \repofile{docs/thesis-report/01-the-pivot/2026-06-26-vision-no-db.md}
  \item Risk grilling with decisions:
    \repofile{findings/no_db_universal_skill/2026-06-26-concept-and-risks.md}
  \item Scope and ordering rationale: \repofile{MASTERPLAN.md} §1--4
\end{itemize}

\subsection*{So what (draft)}
\begin{itemize}
  \item Chosen because it is the only direction that fits every constraint the
    refined problem imposed --- everything database-shaped fails on maintenance,
    everything app-shaped fails on adoption.
\end{itemize}

\subsection*{Now what (draft)}
\begin{itemize}
  \item The proposal directly produced the build plan: prove the university arm
    against a ground truth and a plain-Claude baseline before extending anywhere
    else.
\end{itemize}

\begin{openbox}
\begin{itemize}
  \item Agree on the one-sentence value proposition --- the assignment asks for
    it verbatim and the checklist checks for it.
  \item Were any alternatives seriously discussed beyond app-with-DB vs.\
    skill-without-DB (e.g.\ keeping the app but dropping the DB, a browser
    extension, plain prompt templates)? If yes, name them and why they lost; if
    no, say the decision space was really those two.
\end{itemize}
\end{openbox}

% ---------------------------------------------------------------------------
\section{Building}

The product itself is submitted separately in the Product Artifact ZIP --- per
the assignment, this section only refers to it.

One connecting sentence (draft): the build followed the priority order from
Section~4 --- university discovery first, then a company arm on the same
no-database principle, then consolidation into a single entry-point skill; the
ZIP contains the skill package, install guide, and evaluation harness.

% ---------------------------------------------------------------------------
\section{Testing}

\asks{what was tested, who tested it, how, main feedback, and problems /
surprises / confirmations --- with test notes and quotes as evidence.}

\subsection*{What we have}
\begin{itemize}
  \item Extensive \textbf{self-run evaluation} (this is real testing and should
    be presented as such, but it is not user testing):
    \begin{itemize}
      \item Recall/precision measured per faculty against hand-curated ground
        truth, with the skill scoring far above a plain-Claude baseline.
      \item A steering test showing that different student profiles genuinely
        produce different option maps (not one generic answer).
      \item An independent, deliberately skeptical review that challenged our
        own ``it works'' verdict and produced a punch list --- including
        catching that some of our numbers weren't as clean as claimed. Honest
        material, good for the portfolio.
      \item A surprise worth telling: our first eval design was circular (the
        fixture setup let the system be graded on data it had effectively seen)
        --- we caught it and moved to live evals. That's a genuine ``problem
        discovered through testing.''
    \end{itemize}
  \item \textbf{Informal in-between tests} with people around the project
    happened, including a hands-on session with the course professor that
    produced direct feedback --- but these were never written down properly.
  \item \textbf{Real student beta test: not yet run.} Protocol, capture sheet,
    and recruiting message are written and ready; nobody outside the project has
    used the tool independently yet. The portfolio must state this honestly.
\end{itemize}

\subsection*{Evidence to attach}
\begin{itemize}
  \item Evaluation summary and numbers:
    \repofile{docs/thesis-report/03-hardening-and-evaluation/README.md},
    \repofile{findings/no_db_universal_skill/2026-07-03-eval-aggregate-scorecard.md}
  \item The one clean measurement on the current architecture:
    \repofile{findings/no_db_universal_skill/2026-08-08-theology-blind-run.md}
  \item Skeptical review:
    \repofile{findings/no_db_universal_skill/2026-07-05-fable-1.0-readiness-review.md}
  \item Beta-test protocol (ready, unexecuted):
    \repofile{findings/no_db_universal_skill/2026-08-08-beta-test-protocol.md}
\end{itemize}

\subsection*{So what / Now what}
\begin{itemize}
  \item Depends on the open items below --- this section can't be finished until
    the informal tests are reconstructed and (ideally) the beta test has run.
\end{itemize}

\begin{openbox}
\begin{itemize}
  \item Reconstruct each informal test as a proper record: who tested (role is
    enough, no names needed), what they were asked to do, what they actually
    said --- ideally one verbatim quote each --- and what surprised us or
    confirmed our assumptions.
  \item The professor's hands-on session produced the single most important
    piece of feedback in the whole project (value = reflection, not proposals).
    Reconstruct that session concretely: what was demoed, what was said, what
    changed because of it.
  \item Decide: do we run the beta test before submission? If yes --- results go
    here. If no --- the portfolio says the protocol is ready, explains why it
    hasn't run, and Section~7 carries it as the next step. Either is defensible;
    silently omitting it is not.
\end{itemize}
\end{openbox}

% ---------------------------------------------------------------------------
\section{Iteration and Next Steps}

\asks{changes made from feedback, changes planned but not done, remaining
limitations, distribution/shipping strategy, maintenance/handover strategy, and
the possible role of AI in maintaining or evolving the project.}

\subsection*{What we have}
\begin{itemize}
  \item Changes made based on feedback (pick the 2--3 clearest for the
    portfolio):
    \begin{itemize}
      \item The pivot itself --- the largest feedback-driven change in the
        project.
      \item The professor's reframing (``the value is the student's
        reflection'') was built into the product: the skill now records the
        student's thesis self-understanding before and after the search.
      \item Multiple concrete skill fixes that came directly out of eval
        findings (e.g.\ verification rules that were losing correct candidates).
    \end{itemize}
  \item Changes planned but not implemented: the real student beta test;
    re-measuring the faculties whose numbers predate the architecture change; a
    designed-but-never-run experiment testing our central claim (that the
    skill's advantage comes from curated local scope and would erode as scope
    grows).
  \item Remaining limitations: only one clean measurement on the current
    architecture; company arm weaker than the university arm; coverage honesty
    is a feature but also a limit.
  \item Distribution/shipping (documented, real): the skill package is publicly
    released on GitHub with an install guide, and a cold install was tested.
    Distribution = ``install these markdown folders,'' which is the whole point.
  \item Maintenance/handover: a survival plan exists --- because the product is
    markdown-only (no database, no backend, no API keys), almost nothing can
    rot; main risks are the skill format changing, university sites
    restructuring, and rules going stale. Estimated cost: a couple of hours per
    semester. Ownership options are ranked, \textbf{but no owner has been
    named.}
  \item Role of AI in maintenance --- we have an unusually strong answer: the
    repo ships an operating guide written \emph{for future AI agents}
    (\repofile{AGENTS.md}), so the intended maintainer of this AI-built product
    is itself an AI agent, with a human owner spending $\sim$2\,h/semester
    steering it. This maps exactly onto the assignment's question and should be
    featured, not buried.
\end{itemize}

\subsection*{Evidence to attach}
\begin{itemize}
  \item Open work and rejected ideas:
    \repofile{docs/thesis-report/04-open-work/README.md}
  \item Survival/maintenance plan:
    \repofile{findings/no_db_universal_skill/2026-08-09-survival-and-maintenance-plan.md}
  \item Public release + install guide: GitHub release page,
    \repofile{INSTALL.md}
  \item Agent maintainer guide: \repofile{AGENTS.md}
\end{itemize}

\subsection*{So what (draft)}
\begin{itemize}
  \item Candidate for the most important iteration learning: every major
    improvement came from an outside signal (interviews, professor feedback, the
    skeptical review) --- never from us admiring our own output. To be
    confirmed/sharpened by the group.
\end{itemize}

\subsection*{Now what (draft)}
\begin{itemize}
  \item If the project continued, the next meaningful step is putting the tool
    in real students' hands (the beta test) --- because it's the only claim in
    the whole journey we haven't tested yet.
\end{itemize}

\begin{openbox}
\begin{itemize}
  \item Ownership must be resolved for the portfolio: does one of us take the
    $\sim$2\,h/semester maintenance role, does it go to someone at the
    university, or do we honestly write ``unowned --- and here is what that
    means for the product's lifespan''? Any of the three works; the gap can't
    stay silent.
  \item Has anything changed since the last documented project state (new runs,
    new feedback, new decisions)? The portfolio should reflect where we actually
    are at submission, not where the repo log stops.
  \item Confirm or replace the ``most important iteration learning'' draft above
    --- this should be a sentence the whole group stands behind.
\end{itemize}
\end{openbox}

% ---------------------------------------------------------------------------
\section*{Submission checklist --- current status}
\addcontentsline{toc}{section}{Submission checklist}

\renewcommand{\arraystretch}{1.25}
\noindent\begin{tabularx}{\textwidth}{@{}p{6.2cm}X@{}}
\toprule
\textbf{Checklist item} & \textbf{Status} \\
\midrule
First problem idea & \stdone{} covered (Section~1) \\
Initial project idea / proposal & \stdone{} covered --- \emph{needs one early artifact} \\
Relevant users / user groups & \stdone{} covered --- \emph{needs persona + student-side account} \\
Initial and refined problem definitions & \stpartial{} content there --- \emph{needs the two one-sentence versions} \\
Proposed solution + value proposition & \stpartial{} content there --- \emph{needs the agreed one-sentence value prop} \\
High-level product description & \stdone{} covered (Section~5 + ZIP) \\
User testing and feedback & \stgap{} weakest point --- \emph{informal tests must be written up; beta decision needed} \\
Iteration / planned changes & \stdone{} covered (Section~7) \\
Shipping / access / maintenance strategy & \stdone{} covered --- \emph{except the unresolved owner} \\
Concrete evidence, not just claims & \stdone{} strong --- repo is full of citable artifacts \\
Clear and concise structure & pending --- this draft becomes 4--6 pages once gaps close \\
\bottomrule
\end{tabularx}

\end{document}
